\documentclass[11pt,a4paper]{article}
\usepackage[T1]{fontenc}
\usepackage[utf8]{inputenc}
\usepackage{lmodern}
\usepackage{amsmath,amssymb}
\usepackage[margin=25mm]{geometry}
\usepackage[hidelinks]{hyperref}
\hypersetup{pdftitle={Which limits of Newton's construction reach the zero branch},pdfauthor={navstokgap research working notes}}
\newcommand{\tightlist}{\setlength{\itemsep}{0pt}\setlength{\parskip}{0pt}}
\setlength{\emergencystretch}{3em}
\setcounter{secnumdepth}{0}
\begin{document}
\section{Which limits of Newton's construction reach the zero
branch}\label{which-limits-of-newtons-construction-reach-the-zero-branch}

The thesis under test (user, 2026-10-01): Newton's mechanics exists only
with a positive action scale, no zero branch is continuously reachable,
and the family of theories has a gap in the sense of the mass gap. This
note states the thesis as a claim about limits and proves, in the
simplest cases, which of Newton's own limits reach the zero theory and
which do not. Result: with Gaussian preparations and Gaussian marks,
Galileo's comparison has a record family that is the same for every
value of the constant (Theorem 1), so there is no gap there. For a
determinate preparation, one definite velocity at each place, the
complete coordinate record has no limit as the constant tends to zero
once the preparation has folded, or a limit different from Newton's,
while every smeared record converges to Newton's branch-weighted density
(Theorem 2, any finite partition; Theorem 3, exactly, for inertial
motion with a graded velocity after its fold time). The gap thesis
therefore holds in the topology of complete records of determinate
preparations, beyond their first fold, and fails in the topology of
smeared records or for statistical preparations. Refinement with records
is the one Newtonian limit in which the zero theory is the regular side.
Proofs are written derivations by Claude Fable (2026-10-01), unrefereed;
the stationary phase and almost-periodicity facts used are standard and
named.

Throughout, \(h\) denotes the reduced phase constant, as in the
\href{rivero-1998-conjecture-central-forces.md}{1998-conjecture note};
the recording notes write \(\kappa=h/2\).

\subsection{1. The thesis as a statement about
limits}\label{the-thesis-as-a-statement-about-limits}

For each \(h>0\) let \(T_h\) be the theory with phase constant \(h\):
its dynamics and its record law. The dilation \(h\mapsto\mu h\) is an
isomorphism (\href{principia-fifth-postulate.md}{fifth-postulate note,
Theorem B and §7}; \href{https://arxiv.org/abs/quant-ph/9803035}{Rivero
1998, eq. (10)}, full read), so \(\{T_h:h>0\}\) is one theory in
different units. The zero theory \(T_0\) is Newton's. The thesis says
that \(T_0\) is not the limit of \(T_h\) as \(h\to0\) in any topology
intrinsic to Newton's construction: the scale is generated by the
construction, and zero is isolated. This is the mechanical form of a
mass gap, where the continuum limit generates the scale and the massless
theory is not reached by the coupling.

Known results fix what the thesis can mean. The limit \(h\to0\) is
singular, with phases \(e^{iS/h}\) that have no pointwise limit
{[}@Berry2002; metadata{]}. Yet in weak topologies and for finite times
the classical theory is reached: coherent states follow classical
trajectories for finite times {[}@Hepp1974; metadata{]}, Egorov's
theorem transports smooth observables classically up to \(O(h)\) for
finite times, and Wigner measures give the weak limits of densities
{[}@LionsPaul1993; metadata{]}. So the thesis is false for smeared
records at finite time, and can be true only in stronger topologies.
Newton's construction supplies three candidates: complete records, the
densities themselves rather than their smearings; the all-time limit;
and refinement, with or without records at the inserted times. The rest
of the note tests each.

Deformation families cannot see a gap. The Moyal family, the Fisher
field law and its sheet realizations
(\href{score-constrained-ensemble.md}{score-constrained note}, §§6--8)
are continuous in the constant by construction and contain the zero
branch. The tangent groupoid glues the \(\varepsilon=0\) fibre
continuously to the \(\varepsilon>0\) fibres
(\href{principia-fifth-postulate.md}{fifth-postulate note, §7}), so the
kinematics is continuous at zero as well. A gap, if it exists, lives in
the dynamics and the records.

\subsection{2. Theorem 1: Galileo's comparison has no
gap}\label{theorem-1-galileos-comparison-has-no-gap}

\textbf{Theorem 1 (Gaussian preparations, quadratic forces, Gaussian
marks).} Let \(H=p^2/2m+V\) with \(\deg V\le2\) (inertia, uniform
gravity, Hooke's force). For each admissible \(h\ge0\) prepare the body
with the Gaussian phase-space law \(N(z_0,\Sigma)\), \(\Sigma>0\) fixed;
for \(h>0\) this is the Gaussian state with that Wigner function, which
exists iff \(\det\Sigma\ge h^2/4\). Mark the body at times \(t_j\) by
Gaussian coordinate copies of widths \(\sigma_j\), each read or
forgotten. Then the joint law of all read records and of the final
coordinate is, for every admissible \(h\), the classical Gaussian
process of \(T_0\) with one addition: each mark, read or not, adds to
the body momentum an independent centred Gaussian kick of variance
\(h^2/(4\sigma_j^2)\). The family is continuous in \(h\) and equals
\(T_0\) at \(h=0\).

\emph{Proof.} For \(\deg V\le2\) every Moyal flow is the classical
affine symplectic flow, and the Wigner function is transported
classically (\href{principia-fifth-postulate.md}{fifth-postulate note,
Theorem A(b)}). A coordinate copy is generated by \(q\pi\), quadratic in
the joint canonical variables, so the same holds for the body--pointer
shear. A Gaussian Wigner function of a product state is a positive
density, and reading the pointer coordinate is classical conditioning of
that density. The pointer's momentum \(\pi\) has variance
\(h^2/(4\sigma_j^2)\) and is independent of its coordinate for a minimal
Gaussian pointer, so the shear transfers to the body a kick of that
variance whether or not the coordinate is read; this is Proposition 1 of
the \href{newton-record-parallel-move.md}{parallel-move note} for the
forgotten case, and the read case adds the classical Gaussian
conditioning on \(r=q+z\). Composition in time is the classical flow
between marks. The variance \(h^2/(4\sigma_j^2)\) is continuous in \(h\)
and zero at \(h=0\). \(\square\)

So on the comparison as the Planck paper sets it, the record family is
constant in \(h\) up to continuous kicks. The zero branch is reached in
every topology: there is no gap in Galileo's parabola with statistical
preparations. The same holds for refinement with records at every
inserted time: the fixed-\(h\) limit exists iff
\(\sum_j\sigma_j^{-2}<\infty\)
(\href{newton-record-parallel-move.md}{parallel-move note}), the \(h=0\)
limit always exists, and the \(h\to0\) limit at any fixed partition is
continuous. In this limit the zero theory is the regular side and
\(h>0\) the singular one.

\subsection{3. Theorem 2: determinate preparations with several
paths}\label{theorem-2-determinate-preparations-with-several-paths}

Take the symmetric discrete action of the 1998-conjecture note, (1)
there, on a partition with \(N\) cells of lengths \(\tau_j\), smooth
\(V\), fixed initial point \(q_0=x\), record variable \(y=q_N\), and
internal vertices \(z=(q_1,\ldots,q_{N-1})\in\mathbb R^d\), \(d=N-1\).
With a cutoff \(\chi\in C_c^\infty(\mathbb R^d)\) and a smooth amplitude
\(O\), the halved functional with the record variable kept free is
\[A_h(y)=(2\pi h)^{-d/2}\int\chi(z)O(z)\,e^{iS_\pi(x,z,y)/h}\,dz.\tag{1}\]
Its modulus squared \(|A_h(y)|^2\) is the complete coordinate record at
the final time; \(\int\varphi(y)|A_h(y)|^2dy\) with
\(\varphi\in C_c^\infty\) is a smeared record.

\textbf{Hypothesis (several nondegenerate paths).} On an open interval
\(Y\) of record values, \(S_\pi(x,\cdot,y)\) has exactly \(K\ge2\)
critical points \(z_a(y)\) in \(\operatorname{supp}\chi\), all
nondegenerate, with \(\chi=1\) near each. By the implicit function
theorem they depend smoothly on \(y\). Write
\[S_a(y)=S_\pi(x,z_a(y),y),\qquad
c_a(y)=\frac{O(z_a(y))\,e^{i\pi\sigma_a/4}}{|\det S_\pi''(z_a(y))|^{1/2}},\tag{2}\]
with \(\sigma_a\) the Hessian signature.

\textbf{Lemma (final momenta are distinct).} \(p_a(y):=S_a'(y)\) equals
\(m(y-q^{(a)}_{N-1}(y))/\tau_{N-1}-\tfrac12\tau_{N-1}V'(y)\), and
\(p_a(y)\neq p_b(y)\) for \(a\neq b\).

\emph{Proof.} The envelope formula gives \(S_a'(y)=\partial_yS_\pi\) at
the critical point, which is the displayed expression. The discrete
Euler--Lagrange equation at vertex \(j\),
\(m(q_{j+1}-q_j)/\tau_j-m(q_j-q_{j-1})/\tau_{j-1}
+\tfrac12(\tau_{j-1}+\tau_j)V'(q_j)=0\), is linear in \(q_{j-1}\) with
coefficient \(m/\tau_{j-1}\ne0\). So a critical point is determined by
\((q_{N-1},q_N)\) by backward recursion, two distinct critical points
have distinct \(q_{N-1}\), and the momenta differ. \(\square\)

\textbf{Theorem 2.} Under the hypothesis, uniformly for \(y\) in compact
subsets of \(Y\), \[|A_h(y)|^2=\sum_a|c_a(y)|^2
+\sum_{a<b}2|c_a(y)c_b(y)|\cos\!\Big(\frac{S_a(y)-S_b(y)}{h}+\theta_{ab}(y)\Big)
+O(h).\tag{3}\] (i) \emph{Complete record.} Fix \(y\in Y\). The limit of
\(|A_h(y)|^2\) as \(h\downarrow0\) exists only if, for every nonzero
value \(\omega\) of \(S_a(y)-S_b(y)\), the amplitudes
\(|c_ac_b|e^{i\theta_{ab}}\) of the pairs with that action difference
sum to zero. For \(K=2\) with \(S_1(y)\neq S_2(y)\) this means
\(c_1(y)c_2(y)=0\). When the limit exists it equals \(\sum_a|c_a|^2\)
plus the cross terms of pairs with equal actions, which is Newton's
branch-weighted density
\(D_{0,\pi}(|O|^2)=\sum_a|O(z_a)|^2/|\det S''_\pi(z_a)|\) only if those
cross terms vanish too. (ii) \emph{Smeared record.} For every
\(\varphi\in C_c^\infty(Y)\),
\[\int\varphi(y)|A_h(y)|^2dy\longrightarrow
\sum_a\int\varphi(y)\frac{|O(z_a(y))|^2}{|\det S''_\pi(z_a(y))|}dy
\qquad(h\downarrow0).\tag{4}\]

\emph{Proof.} Equation (3) is the stationary phase expansion of (1) with
nondegenerate critical points, (3) of the 1998-conjecture note, whose
remainder is uniform in the parameter \(y\) on compact sets where the
critical points stay nondegenerate (Hörmander's stationary phase theorem
with parameters; standard). Expanding the modulus squared gives (3) with
\(\theta_{ab}=\pi(\sigma_a-\sigma_b)/4+\arg O_a-\arg O_b\).

\begin{enumerate}
\def\labelenumi{(\roman{enumi})}
\item
  Put \(s=1/h\). The cross sum in (3) is an almost periodic function
  \(f(s)\) in Bohr's sense, a finite trigonometric sum. If
  \(|A_h(y)|^2\) converges as \(h\downarrow0\) then \(f(s)\) converges
  as \(s\to\infty\). An almost periodic function with a limit at
  infinity is constant: for any \(s_0\) and \(\delta>0\) the set of
  \(\delta\)-almost periods is relatively dense, so there are
  \(T_n\to\infty\) with \(|f(s_0+T_n)-f(s_0)|<\delta\), and the limit
  forces \(f(s_0)=\lim f\). A constant trigonometric sum has zero
  coefficient at every nonzero frequency: the Bohr mean of
  \(f(s)e^{-i\omega s}\) is that coefficient. This proves the condition,
  and the value of the limit when it exists.
\item
  In the cross term with indices \(a\neq b\) the phase
  \(\psi_{ab}=S_a-S_b\) has \(\psi_{ab}'=p_a-p_b\neq0\) on \(Y\) by the
  Lemma. One integration by parts,
  \(\int g\,e^{i\psi/h}dy=ih\int (g/\psi')'\,e^{i\psi/h}dy\) for
  \(g\in C_c^\infty(Y)\), bounds the term by \(O(h)\). The diagonal
  terms and the uniform remainder give (4). \(\square\)
\end{enumerate}

\textbf{Corollary (finite resolution).} Let \(\varphi\ge0\) be smooth
with compact support and \(\int\varphi=1\), and let
\(\varphi_\delta(y)=\delta^{-1}\varphi((y-y_0)/\delta)\) be a record of
resolution \(\delta\) at \(y_0\in Y\). With \(\Delta p>0\) the minimum
of \(|p_a-p_b|\) over pairs \(a\ne b\) on the support,
\[\int\varphi_\delta|A_h|^2dy=\sum_a\int\varphi_\delta
\frac{|O(z_a)|^2}{|\det S''_\pi(z_a)|}dy
+O\Big(\frac{h}{\delta\,\Delta p}\Big)+O(h).\tag{4'}\] A record of
resolution \(\delta\) reaches Newton's density iff
\(h/(\delta\Delta p)\to0\): the fringes, of spacing \(h/\Delta p\), must
be unresolved. The complete-record topology of (i) is the order of
limits in which \(\delta\to0\) before \(h\to0\); the smeared topology of
(ii) is the opposite order.

\emph{Proof.} In the integration by parts of (ii), the derivative of
\(\varphi_\delta\) carries a factor \(\delta^{-1}\) relative to
\(\varphi_\delta\) itself, and \(\int\varphi_\delta=1\); the remaining
terms are as before. \(\square\)

Thus the complete record of a determinate preparation with several paths
does not reach Newton's density as \(h\downarrow0\), by non-convergence
in the generic case of distinct actions and by a wrong limit in the
symmetric case of equal actions, while its smearing always does. Theorem
1 of the 1998-conjecture note is the case \(d=1\), \(K=3\), \(y=0\) of
(i), including the retained cross term of the two mirror paths; smearing
in \(y\) removes that term by (ii), because the mirror paths arrive at
\(y=0\) with opposite momenta.

\subsection{3b. The gap as a number: fringe
visibility}\label{b.-the-gap-as-a-number-fringe-visibility}

A mass gap is a number that is positive for every value of the coupling
and zero for the free theory. The dilation \(h\mapsto\mu h\) forbids any
scale from playing that role here, so the invariant must be
dimensionless. It is the contrast of the fringes in (3).

\textbf{Proposition (visibility invariant).} Let \(y_0\in Y\) be a
record point with two branches, \(K=2\), and write
\(\rho_a=|c_a(y_0)|^2=|O(z_a)|^2/|\det S''_\pi(z_a)|\) for Newton's
branch densities there, \(\Delta p=|p_1-p_2|(y_0)\). (i) For every
\(h>0\) the complete record near \(y_0\) oscillates with fringe spacing
\(2\pi h/\Delta p+O(h^2)\) and visibility
\[V=\frac{\max|A_h|^2-\min|A_h|^2}{\max|A_h|^2+\min|A_h|^2}
=\frac{2\sqrt{\rho_1\rho_2}}{\rho_1+\rho_2}+O(h),\tag{V1}\] the maximum
and minimum taken over one fringe. The leading term is fixed by Newton's
branch densities and is independent of \(h\). At \(h=0\) the record is
\(\rho_1+\rho_2\) and the visibility is zero. (ii) The limit set of
\(|A_h(y_0)|^2\) as \(h\downarrow0\) is the interval
\((\rho_1+\rho_2)\,[1-V,\,1+V]\). (iii) For a record of resolution
\(\delta\) with kernel \(\varphi\), and \(\delta\) small enough that the
fringe phase is affine across the window (\(\delta^2|p_1'-p_2'|\ll h\)),
the visibility is
\[V_\delta=V\,|\hat\varphi(\delta\Delta p/h)|+O(h),\qquad
\hat\varphi(\xi)=\int\varphi(u)e^{i\xi u}du,\tag{V2}\] so a Gaussian
kernel gives \(V_\delta=V\exp[-(\delta\Delta p/h)^2/2]\). The contrast
is seen iff \(h\gg\delta\Delta p\), and the two orders of limits of the
Corollary give \(V\) and \(0\).

\emph{Proof.} (i) By (3) with \(K=2\),
\(|A_h(y)|^2=\rho_1+\rho_2+2\sqrt{\rho_1\rho_2}\cos(\psi(y)/h+\theta)+O(h)\)
with \(\psi=S_1-S_2\) and \(\psi'=p_1-p_2\neq0\), uniformly near
\(y_0\). Over one fringe the cosine runs through \(\pm1\), which gives
the spacing, the extremes \(\rho_1+\rho_2\pm2\sqrt{\rho_1\rho_2}\) and
(V1). (ii) \(\psi(y_0)\neq0\) generically, and \(\psi(y_0)/h\) is
continuous and unbounded in \(h\), so its residue mod \(2\pi\) takes
every value along sequences \(h\downarrow0\); the cosine takes every
value in \([-1,1]\). If \(\psi(y_0)=0\) the cross term is constant and
the limit is the top of the interval. (iii) Integrate the cross term
against \(\varphi_\delta\): with
\(\psi(y)=\psi(y_0)+\Delta p\,(y-y_0)+O(\delta^2|\psi''|)\) on the
window, the integral is
\(2\sqrt{\rho_1\rho_2}\,\mathrm{Re}[e^{i\psi(y_0)/h+i\theta}\hat\varphi(\delta\Delta p/h)]\)
up to the stated error; the diagonal terms are unchanged. \(\square\)

With \(K\ge3\) branches, as after a fold, each pair carries its own
contrast \(2\sqrt{\rho_a\rho_b}/\sum_c\rho_c\) at its own spacing
\(2\pi h/|p_a-p_b|\). For inertia with graded velocity (Theorem 3) the
branch densities are \(a(q_a)^2/|1+p_0'(q_a)t/m|\); for the Kepler swarm
(Theorem 4) they are \(a(\phi'_a)^2/|1+I_0'\Omega't|_a\). In every case
the invariant is computed from Newton's multi-stream density and is
positive wherever two streams cross, for every \(h>0\), and zero at
\(h=0\). That is the gap: a contrast rather than a scale, discontinuous
at zero, and invisible at any fixed resolution once \(h\) falls below
\(\delta\Delta p\).

\subsection{3c. Imprecision actions at both
ends}\label{c.-imprecision-actions-at-both-ends}

The record is one end of the comparison; the preparation is the other. A
determinate swarm attached to the \(h>0\) theory as one state (5) has no
momentum thickness; a real swarm has one, and so does the same classical
swarm attached as a mixture of localized states, whose thickness is of
order \(\sqrt h\). The following proposition says what a thickness does,
in the exact free setting of §4, and the general case follows by the
same envelope identity.

\textbf{Proposition (preparation thickness).} Shift the initial curve by
a momentum \(\eta\), \(\psi_0^\eta=a\,e^{i(S_0+\eta q)/h}\), and average
the record over \(\eta\) with a density \(w\) of mean zero and width
\(\Delta\). Let \(x\) be a record point with two arrivals
\(q_1\neq q_2\), write \(\Delta q_0=|q_1-q_2|\) for their launch
separation, and let \(\Delta\) be small enough that the branch densities
and launch points are constant across the thickness up to \(o(1)\) and
\(\Delta^2|\partial_\eta(q_1-q_2)|\ll h\). Then the averaged record has
visibility \[V_\Delta=V\,|\hat w(\Delta q_0/h)|+O(h)+o(1),\qquad
\hat w(\xi)=\int w(\eta)e^{i\xi\eta}d\eta,\tag{V3}\] so a Gaussian
thickness gives \(V_\Delta=V\exp[-(\Delta\,\Delta q_0/h)^2/2]\).

\emph{Proof.} The arrivals solve \(q+(p_0(q)+\eta)t/m=x\) and the
stationary actions are \(\Phi^\eta_a=m(x-q_a)^2/2t+S_0(q_a)+\eta q_a\).
By the envelope identity, \(d\Phi^\eta_a/d\eta=q_a(\eta)\), so
\(\Phi^\eta_1-\Phi^\eta_2=\psi(0)+\eta(q_1-q_2)+O(\eta^2|\partial_\eta(q_1-q_2)|)\).
Averaging the cross term of (3) against \(w\) under the stated
hypotheses gives
\(2\sqrt{\rho_1\rho_2}\,\mathrm{Re}[e^{i\psi(0)/h+i\theta}\hat w((q_1-q_2)/h)]\),
and the diagonal terms change by \(o(1)\). \(\square\)

With both imprecisions present and independent, the contrast is
\(V\,|\hat\varphi(\delta\Delta p/h)|\,|\hat w(\Delta\,\Delta q_0/h)|\).
Each argument is an action divided by \(h\): the record imprecision
\(\delta\,\Delta p\), position resolution times the arrival momentum
difference, and the preparation imprecision \(\Delta\,\Delta q_0\),
momentum thickness times the launch separation. The gap is seen iff both
imprecision actions lie below \(h\). This is the floor the thesis
asserts, and it runs opposite to the uncertainty relation: there \(h\)
bounds imprecisions from below; here imprecisions must be pushed below
\(h\) for the \(h>0\) record to differ from Newton's. Newton's
construction has exact preparations and exact records, both imprecision
actions zero, so for every \(h>0\) it sees the full contrast \(V\), and
at \(h=0\) it sees none. The two attachments of the same classical swarm
are now separated by a number: the mixture of localized states has
\(\Delta\sim\sqrt h\), so \(\Delta\,\Delta q_0/h\to\infty\) and its
contrast vanishes, while the state (5) has \(\Delta=0\). Determinacy at
each \(h\), momentum thickness of order \(h\) at fixed position (the
sheet separation \(2\kappa\,\partial_q\log\rho\) of the
\href{score-constrained-ensemble.md}{score note}) rather than
\(\sqrt h\), is the premise that selects the first attachment.

\subsection{4. Theorem 3: inertial motion with a graded
velocity}\label{theorem-3-inertial-motion-with-a-graded-velocity}

Newton's simplest motion already shows the gap when the preparation is
determinate and its velocity varies along it. Let \(h>0\), \(V=0\), and
prepare
\[\psi_0(q)=a(q)e^{iS_0(q)/h},\qquad a\in C_c^\infty(\mathbb R)\ \text{real},
\qquad p_0=S_0'.\tag{5}\] Each particle at \(q\) has the definite
velocity \(p_0(q)/m\). The free propagator is exact, so
\[\psi_t(x)=\Big(\frac{m}{2\pi iht}\Big)^{1/2}\int a(q)\,
e^{i\Phi_x(q)/h}dq,\qquad
\Phi_x(q)=\frac{m(x-q)^2}{2t}+S_0(q).\tag{6}\] Critical points of
\(\Phi_x\) are the arrivals \(q+p_0(q)t/m=x\), and
\(\Phi_x''(q)=m/t+p_0'(q)\). Define the fold time
\[t_*=\frac{m}{\max(-p_0')}\in(0,\infty],\tag{7}\] with \(t_*=\infty\)
when \(p_0'\ge0\) everywhere.

\textbf{Theorem 3.} (a) For \(0<t<t_*\) the arrival map
\(q\mapsto q+p_0(q)t/m\) is a diffeomorphism, every \(x\) has one
arrival \(q_*(x)\), and
\(|\psi_t(x)|^2\to a(q_*(x))^2/(1+p_0'(q_*(x))t/m)\) pointwise, the
classical transported density: the complete record reaches the zero
branch before the fold. (b) For \(t>t_*\) there is an open interval of
\(x\) with three nondegenerate arrivals \(q_1<q_2<q_3\). Where \(a\) is
nonzero at two or more of them, Theorem 2 applies with \(d=1\), \(K=3\):
\(|\psi_t(x)|^2\) has no limit as \(h\downarrow0\) unless the cross
amplitudes cancel, and every smearing in \(x\) converges to the
classical three-branch density \(\sum_a a(q_a)^2/|1+p_0'(q_a)t/m|\). (c)
At a fold point \(x_c(t)\), where \(\Phi_x''(q_c)=0\) and
\(\Phi_x'''(q_c)\neq0\), \(|\psi_t(x_c)|^2\) grows like \(h^{-1/3}\):
the record diverges with the classical caustic density.

\emph{Proof.} (a) The derivative \(1+p_0'(q)t/m\) is positive for
\(t<t_*\), so the arrival map is strictly increasing and onto, and the
single critical point is nondegenerate; one-dimensional stationary phase
in (6) gives \(|\psi_t|^2\to a^2/|\Phi_x''|\cdot(m/t)\), which is the
stated ratio. (b) For \(t>t_*\) the derivative is negative on an open
set, so the arrival map is not monotone; by the intermediate value
theorem the preimage of each \(x\) in an open interval has three points,
nondegenerate away from the two fold values. The amplitudes
\(c_a=a(q_a)e^{i\pi\sigma_a/4}/|\Phi_x''(q_a)|^{1/2}\) and the actions
\(\Phi_x(q_a)\) play the roles of (2); the final momenta
\(\partial_x\Phi_x(q_a)=m(x-q_a)/t=p_0(q_a)\) are distinct because the
\(q_a\) are distinct and \(x=q_a+p_0(q_a)t/m\). Theorem 2 applies. (c)
At a degenerate critical point of order three the stationary phase
integral scales as \(h^{1/3}\) against the prefactor \(h^{-1/2}\),
giving \(|\psi|^2\sim h^{-1/3}\) (Airy scaling; standard). \(\square\)

Newton's inertial motion, with a preparation in which faster particles
start behind slower ones, folds at \(t_*\). Before the fold the zero
branch is reached by complete records; after it, never. The fringe
spacing in \(x\) is \(h/|p_a-p_b|\), which tends to zero without the
record ever settling. The classical multi-stream density is recovered
only by smearing, an operation external to the determinate construction.

\subsection{4b. Theorem 4: central forces on the radial cycle, Hooke
against
Kepler}\label{b.-theorem-4-central-forces-on-the-radial-cycle-hooke-against-kepler}

Newton's central forces supply the case the thesis was made for. Fix the
angular momentum \(\ell>0\). The radial motion is one degree of freedom
with action \(I=I_r\) and angle \(\phi\) on the circle, Hamiltonian
\(E(I)\), frequency \(\Omega(I)=E'(I)\) and radial period
\(T=2\pi/\Omega\). The
\href{rivero-1998-conjecture-central-forces.md}{1998-conjecture note,
(10)--(11)} gives, for Hooke's force (Book I, Proposition X),
\(E=\omega(2I+\ell)\), so \(\Omega=2\omega\) and \(\Omega'=0\), the
isochrony; and for Kepler's, \[E=-\frac{mk^2}{2(I+\ell)^2},\qquad
\Omega=\frac{mk^2}{(I+\ell)^3},\qquad
\Omega'=-\frac{3mk^2}{(I+\ell)^4},\qquad
T=2\pi\sqrt{\frac{ma^3}{k}}.\tag{8}\] For Kepler the angle \(\phi\) is
the mean anomaly, which Kepler's equation turns into position on the
ellipse: it is the record astronomy makes.

\emph{Model.} Take \(H_h=E(-ih\partial_\phi)\) on the circle with the
twisted periodicity \(\psi(\phi+2\pi)=e^{2\pi i\alpha}\psi(\phi)\),
\(\alpha\) the Maslov shift of the radial cycle (\(\alpha=\tfrac12\)),
so that the spectrum is \(E(h(n+\alpha))\), the Keller--Maslov values,
which for Hooke and Kepler are the exact planar spectra ((9)--(11) of
that note). Extend \(E\) smoothly outside the physical range of \(I\);
for the preparations below this changes the evolved state by
\(O(h^\infty)\). The determinate preparation is a swarm,
\[\psi_0(\phi)=a(\phi)e^{iS_0(\phi)/h},\qquad
a\in C_c^\infty((-\pi,\pi)),\qquad I_0=S_0',\tag{9}\] each body at
\(\phi\) on the orbit of action \(I_0(\phi)\), with action spread
\(\Delta I=\max I_0-\min I_0\) on the support and support length
\(|a|<2\pi\). Poisson summation over the eigenvalue index gives exactly
\[\psi_t(\phi)=\sum_{k\in\mathbb Z}\frac{e^{-2\pi ik\alpha}}{2\pi h}
\iint a(\phi')\,e^{i\Phi_k(\phi',I)/h}\,d\phi'\,dI,\qquad
\Phi_k=S_0(\phi')+I(\phi-\phi'+2\pi k)-E(I)t.\tag{10}\] The critical
points of \(\Phi_k\) are \(I=I_0(\phi')\) and
\(\phi-\phi'+2\pi k=\Omega(I)t\): the body prepared at \(\phi'\) arrives
at \(\phi\) after \(k\) turns. The Hessian determinant is
\(-(1+I_0'(\phi')\Omega'(I_0(\phi'))t)\). Define the fold time and the
lapping time \[t_*=\frac1{\max(-I_0'\,\Omega'(I_0))}\in(0,\infty],\qquad
t_{\rm lap}=\frac{2\pi-|a|}{\max\Omega(I_0)-\min\Omega(I_0)}\in(0,\infty].\tag{11}\]

\textbf{Theorem 4.} (a) \emph{Hooke.} \(E\) is affine in \(I\), the
\(I\)-integral in (10) is a delta function, and for every \(h>0\) and
every \(t\)
\[|\psi_t(\phi)|^2=a(\phi-2\omega t)^2\quad(\text{mod }2\pi):\tag{12}\]
the swarm rotates rigidly, the complete anomaly record is the same for
every \(h\), and the zero branch is reached at all times. (b)
\emph{Kepler, or any \(\Omega'\neq0\).} For \(t<\min(t_*,t_{\rm lap})\)
every \(\phi\) has at most one arrival and
\(|\psi_t(\phi)|^2\to a(\phi')^2/|1+I_0'(\phi')\Omega'(I_0(\phi'))t|\)
pointwise. For \(t>t_{\rm lap}\) some \(\phi\) have two or more arrivals
of different winding, and for \(t>t_*\) three arrivals of the same
winding; distinct arrivals at the same \(\phi\) have distinct actions
\(I_a\), so Theorem 2 and its Corollary apply with \(\phi\) as record
variable and \(I\) as momentum: the complete anomaly record has no
\(h\to0\) limit unless cross amplitudes cancel, every smeared record
converges to the classical density
\(\sum_a a(\phi'_a)^2/|1+I_0'\Omega't|_a\), and a record of angular
resolution \(\delta\) reaches it iff \(h/(\delta\min|I_a-I_b|)\to0\).
(c) \emph{Kepler times.} For a narrow swarm, \(\Delta I\ll I+\ell\),
\[t_{\rm lap}\simeq\frac{2\pi(I+\ell)^4}{3mk^2\,\Delta I}
=\frac{T\,(I+\ell)}{3\,\Delta I}=\frac{T^2}{\Delta T},\qquad
t_*=\frac{(I+\ell)^4}{3mk^2\max I_0'}\ \ (\max I_0'>0),\tag{13}\] with
\(\Delta T=3T\Delta I/(I+\ell)\) the spread of periods. Both are
classical times, independent of \(h\).

\emph{Proof.} The eigenfunctions are \(e^{i(n+\alpha)\phi}\) with
eigenvalues \(E(h(n+\alpha))\), and the preparation's coefficients are
\(c_n=(2\pi)^{-1}\int a\,e^{i[S_0(\phi')-h(n+\alpha)\phi']/h}d\phi'\).
With \(\nu=n+\alpha\) and \(I=h\nu\) the evolved state is
\(\sum_nG(n+\alpha)\) for
\(G(\nu)=(2\pi)^{-1}\int a(\phi')e^{i[S_0(\phi')+h\nu(\phi-\phi')-E(h\nu)t]/h}d\phi'\),
a Schwartz function of \(\nu\): for \(h\nu\) outside the range of
\(I_0\) the \(\phi'\)-phase has no stationary point and integration by
parts gives rapid decay, and derivatives in \(\nu\) only add polynomial
factors. Poisson summation,
\(\sum_nG(n+\alpha)=\sum_ke^{2\pi ik\alpha}\int G(\nu)e^{-2\pi ik\nu}d\nu\),
with \(k\to-k\) and \(d\nu=dI/h\), is (10). (a) For
\(E=\omega(2I+\ell)\) the \(I\)-integral is
\((2\pi h)^{-1}\int e^{iI(\phi-\phi'+2\pi k-2\omega t)/h}dI
=\delta(\phi-\phi'+2\pi k-2\omega t)\), so
\(\psi_t(\phi)=\sum_ke^{-2\pi ik\alpha-i\omega\ell t/h}
a(\phi+2\pi k-2\omega t)e^{iS_0(\phi+2\pi k-2\omega t)/h}\); one \(k\)
contributes at each \(\phi\) because \(|a|<2\pi\), and (12) follows. (b)
The arrival map \(\phi'\mapsto\phi'+\Omega(I_0(\phi'))t\) has derivative
\(1+I_0'\Omega't>0\) for \(t<t_*\), so it is injective on the support,
and its image has length \(|a|+t(\max\Omega-\min\Omega)<2\pi\) for
\(t<t_{\rm lap}\), so no two windings overlap: one nondegenerate
critical point at each \(\phi\) of the image, none elsewhere up to
\(O(h^\infty)\). Two-dimensional stationary phase in (10) with the
prefactor \((2\pi h)^{-1}\) gives the pointwise limit. For
\(t>t_{\rm lap}\) the image length exceeds \(2\pi\) and by continuity
some \(\phi\) is covered by two points of the support with different
\(k\); for \(t>t_*\) the map is not monotone and three arrivals of the
same winding appear, as in Theorem 3. Two arrivals \((\phi'_1,k_1)\),
\((\phi'_2,k_2)\) at the same \(\phi\) with the same action \(I\) would
satisfy \(\phi'_1-\phi'_2=2\pi(k_1-k_2)\), hence coincide because
\(|a|<2\pi\); so distinct arrivals have distinct actions, and
\(\partial_\phi\Phi_k=I\) at the critical point is the record momentum
of Theorem 2. (c) Insert (8) into (11): \(|\Omega'|\Delta I\) is the
frequency spread, \(T=2\pi(I+\ell)^3/(mk^2)\), and
\(T\propto(I+\ell)^3\) gives \(\Delta T/T=3\Delta I/(I+\ell)\).
\(\square\)

Proposition X and Kepler's law thus fall on opposite sides of the thesis
in the record Newton's astronomy actually makes. The isochronous force
keeps the zero branch for all time; the Kepler force loses it for every
swarm with a spread of periods, after the classical lapping time
\(T^2/\Delta T\), the time after which the fast bodies have gained a
full turn on the slow ones. From then on the complete record of
anomalies carries fringes of angular spacing \(h/\Delta I\) that never
settle. A stationary preparation, the whole ring at one action, is an
eigenstate and reaches the zero branch trivially; the gap needs a
localized swarm with a spread of actions. In the radial record \(r\), by
contrast, even Hooke's swarm has two branches with opposite radial
momenta whenever it straddles a turning point, so there the gap appears
for both forces during each turning passage; the anomaly record is the
one that separates them.

\subsection{5. All-time invariants}\label{all-time-invariants}

Two quantities already in the corpus show the other non-commuting pair
of limits, \(h\to0\) against \(t\to\infty\). For the free Gaussian, the
expected number of sign reversals of the two-sheet realization over all
time is \(\pi\varphi(1)\) for every \(h>0\) and zero at \(h=0\)
(\href{sed-closure-under-recording.md\#823-a-stationary-switching-process-repairs-the-gaussian-marginal}{recording
note, (150)}); the total phase \(\arctan u\) of (143) there reaches
\(\pi/2\) for every \(h>0\) and is zero at \(h=0\). Both are limits
taken at infinite time first; at fixed finite time both vanish
continuously as \(h\to0\). The first is a hidden count and the second a
phase, so neither is a record; their observable counterpart is the
breakdown of the classical limit after the Ehrenfest time, which Theorem
3 exhibits at the fold time \(t_*\) for inertial motion. The all-time
limit therefore supports the thesis only through folding, which is
already Theorem 3.

\subsection{6. What the thesis is, and what it
needs}\label{what-the-thesis-is-and-what-it-needs}

\begin{itemize}
\tightlist
\item
  \textbf{Galileo's comparison with statistical preparations has no
  gap.} Theorem 1 makes the record family constant in \(h\). Any gap
  argument run on the parabola with Gaussian preparations and marks will
  fail, for the same reason the deformation families fail.
\item
  \textbf{Determinate preparations that fold have the gap, in complete
  records.} Theorems 2 and 3 prove it on every finite partition and
  exactly for inertia with graded velocity; Theorem 4 proves it for
  Kepler's force in the anomaly record. Beyond Galileo the thesis holds
  in this topology.
\item
  \textbf{Hooke and Kepler separate.} In the anomaly record, Proposition
  X's isochronous force keeps the zero branch for all time (Theorem
  4(a)), while Kepler's force loses it for every swarm with a spread of
  periods after the classical lapping time \(T^2/\Delta T\) (Theorem
  4(b), (c)). The gap is a property of the force law.
\item
  \textbf{Smeared records and statistical preparations reach zero.} This
  is the content of the weak-limit theorems named in §1, and of (4).
\item
  \textbf{Refinement with records reverses the roles.} The zero theory
  is the regular side of that limit.
\item
  \textbf{The thesis is an order of limits.} By the Corollary to Theorem
  2, a record of resolution \(\delta\) reaches Newton iff
  \(h/(\delta\Delta p)\to0\). ``Complete records'' means \(\delta\to0\)
  before \(h\to0\): the record limit is taken first, as the continuum
  limit is taken first in the mass gap.
\item
  \textbf{The gap is a contrast, not a scale.} The fringe visibility
  \(2\sqrt{\rho_1\rho_2}/(\rho_1+\rho_2)\) of a complete record where
  two streams cross is fixed by Newton's branch densities, independent
  of \(h\) for every \(h>0\), and zero at \(h=0\) (§3b). The dilation
  symmetry makes the invariant dimensionless where the mass gap is a
  scale.
\item
  \textbf{Both ends must be exact.} The contrast is damped by
  \(|\hat\varphi(\delta\Delta p/h)|\) at the record and by
  \(|\hat w(\Delta\,\Delta q_0/h)|\) at the preparation (§3c): the gap
  is seen iff both imprecision actions lie below \(h\). Newton's
  construction has both zero; a classical swarm attached as a mixture of
  localized states has thickness \(\sqrt h\) and shows no gap.
\end{itemize}

The thesis is therefore exactly as strong as two premises: that records
are complete, the densities themselves rather than their smearings, and
that preparations are determinate, one velocity at each place. Neither
is foreign to Newton. A determinate preparation is the Newtonian
default, every body with its own velocity; the statistical preparation
of Theorem 1 is the extra assumption the Planck paper makes with its
Gaussian marks. The second is the fifth postulate's joint determinacy
(\href{principia-fifth-postulate.md}{fifth-postulate note}) read as a
preparation, Newton's \emph{velocitas ultima} assigned to every particle
of the ensemble. The first is the records reading already required by
\href{leibniz-continuity-records.md}{Leibniz continuity}. It is also
Newton's own topology. His limits are exact: the Section I scholium asks
the reader to think of the vanishing quantities ``semper diminuendas
sine limite'', always diminishing without limit, and the geometry of the
comparison has no floor at any resolution
(\href{planck-gap-paper.md}{Planck paper, §2, Proposition 1}). Newton's
construction contains no apparatus resolution; the ultimate ratio is the
exact record, and the resolution limit is taken inside the geometry
before any constant enters. Read in that topology, the \(h>0\) family is
Theorem 2(i)'s complete record. The smeared topology adds a resolution
\(\delta\) that the construction does not contain; the Planck paper's
marks are that modern addition. The gap claim and the Leibniz claim rest
on the same reading of records, and the gap claim adds determinacy of
the preparation. Under those two premises the statement ``Newton exists
only with a positive scale'' has the precise form: the construction
without branch data exists for every \(h>0\) and refines without extra
data (1998-conjecture note, Theorem 2(b)), while at \(h=0\) it needs a
branch rule (Theorem 1 there) and, after a fold, is not the limit of the
\(h>0\) records (Theorems 2--3 here). With either premise dropped, the
zero branch is reached.

\subsection{7. Consequence for STATE}\label{consequence-for-state}

The gap thesis is located and has its first central-force case. It holds
for complete records of determinate preparations beyond their first fold
(Theorems 2--3), and in the anomaly record it separates the force laws:
Hooke keeps the zero branch for all time, Kepler loses it after
\(T^2/\Delta T\) (Theorem 4). It fails for Galileo's comparison with
Gaussian preparations and marks (Theorem 1), for smeared records, and
for refinement with records; a record of resolution \(\delta\) reaches
Newton iff \(h/(\delta\Delta p)\to0\). The gap has a number: the fringe
visibility \(2\sqrt{\rho_1\rho_2}/(\rho_1+\rho_2)\) of a complete record
where two streams cross, fixed by Newton's branch densities, independent
of \(h\) for every \(h>0\) and zero at \(h=0\) (§3b), seen iff the
record and preparation imprecision actions \(\delta\Delta p\) and
\(\Delta\,\Delta q_0\) both lie below \(h\) (§3c). Consequences: the
Planck paper's Gaussian parabola cannot carry the gap argument, which
belongs to the Kepler swarm; the two premises, complete records as the
record-limit-first order and determinate preparations as Newton's
default, are the gap's hypotheses and the next things to defend, the
first with the Leibniz records reading. Further exact cases available:
the radial record of a Hooke swarm at a turning passage, and the
classical-first route read with \(\varepsilon\) in place of \(h\), where
Theorem 2(i) says the unselected construction at zero exists only for
single-path problems. Reject a gap argument run inside a deformation
family, on Gaussian preparations of the parabola, or on smeared records,
and reject presenting the non-commutation of \(h\to0\) with
\(t\to\infty\) as a record without an observable.
\end{document}
